\chapter{Parachutes, recovery systems, and stage separation}\label{ch:recovery}
Recovery spans three layers: a component describes the device; a deployment rule or FC command decides when deployment is requested; the engine changes the actual flight state. A logical firing pulse, queued event, accepted deployment, and changed descent drag are distinct observations.

\section{Device geometry and mass}
\class{RecoveryDevice} is the common abstraction for parachutes and streamers. It supplies projected area, drag coefficient behavior, material-derived mass, and per-flight deployment configuration. \class{Parachute} derives canopy area from diameter and includes line mass from line count, length, and line material. \class{Streamer} derives area from its dimensions. Packed position and dimensions support placement and mass modeling inside the rocket.

The canopy shown during flight playback is a visual indicator of actual deployment. It is not the mesh used to calculate drag. Likewise, a shock cord component contributes modeled properties but does not create a dynamic tether constraint, elasticity, or opening-shock calculation between separated bodies.

\section{Ordinary OpenRocket deployment rules}
\class{DeploymentConfiguration} is indexed by flight configuration. Its event choices include launch, same-stage motor ejection, apogee, a descending altitude crossing, lower-stage separation, and never. A configured delay is added after the activating event. The engine checks active stages and dispatches the resulting \class{FlightEvent.Type.RECOVERY_DEVICE_DEPLOYMENT} through the listener system.

Altitude deployment detects a crossing using the previous/current altitude pair, not a statement that the rocket is simply below a threshold at any time. Ejection matching is stage-aware. Stage-separation-triggered recovery uses the relevant neighboring-stage relationship. These details prevent rules from being interpreted as interchangeable generic timers.

\section{Inherited parachute and streamer force laws}
A parachute's deployed area is $A=\pi d^2/4$ for the configured canopy diameter. Its baseline drag coefficient is 0.8 unless a preset or explicit setting supplies another value. The current component method returns its coefficient without a Mach-dependent canopy model. Structural mass is canopy area times surface density, plus line count times line length times line density. Canopy porosity, spill holes, inflation history, and suspension-line geometry are not separately resolved in this coefficient.

A streamer's area is strip length times width, and its material mass uses surface density. With automatic CD, the implementation estimates
\begin{equation}
C_{D,\rm streamer}=\min\left[0.4,\ 0.034\frac{\sigma+0.025}{0.105}\frac{\ell+1}{\ell}\right],
\end{equation}
using the source's SI-valued strip length $\ell$ and surface density $\sigma$. The numerical constants belong to this empirical relation. A manual coefficient bypasses the automatic estimate through \class{RecoveryDevice}'s coefficient-selection path. Width still changes drag through area even though it does not appear explicitly in this CD formula.

Once deployed, each device contributes $C_D A$ to the landing stepper. Increasing canopy diameter therefore changes both drag area and material mass; entering an explicit mass override can separate those effects. Multiple deployed devices add their effective drag areas. There is no canopy-shadowing or entanglement calculation between a drogue and main.

For intuition only, balancing weight against drag in still air at fixed density gives $V_t\approx\sqrt{2mg/(\rho\sum C_DA)}$. The simulation integrates acceleration instead of instantly assigning this terminal speed, and wind and altitude-dependent density change the trajectory. This estimate is a useful interpretation of the implemented area sum, not an additional algorithm in the code.

\section{Ordinary staging and two-device recovery}
Independent deployment configurations let a drogue activate at apogee and a main activate on a descending-altitude crossing with its own delay. They share the stage's evolving status but retain separate device identity. A never-deploy rule leaves a device as packed mass unless another authorized path deploys it. The configured deployment threshold is evaluated through altitude events, so it should not be confused with geographic MSL altitude or the FC's zeroed barometric height.

Stage separation has its own event selector: launch; current-stage ignition, burnout, or ejection; upper-stage ignition; ascending altitude; apogee; descending altitude; or never. A separate delay applies after activation. A separation trigger does not itself prescribe a parachute opening, and a motor ejection event does not necessarily separate a stage unless the stage rule says so. After a split, active-component filtering determines the motors and recovery devices available to each branch.

\sourceline{\source{core/src/main/java/info/openrocket/core/rocketcomponent/Parachute.java}{Parachute}; \source{core/src/main/java/info/openrocket/core/rocketcomponent/Streamer.java}{Streamer}; \source{core/src/main/java/info/openrocket/core/rocketcomponent/RecoveryDevice.java}{RecoveryDevice coefficient selection}; \source{core/src/main/java/info/openrocket/core/rocketcomponent/StageSeparationConfiguration.java}{StageSeparationConfiguration}.}


\section{Exclusive FC recovery mode}
When FC simulation is enabled, both event engines suppress ordinary automatic deployment scheduling. Only events owned by the FC recovery bridge are accepted as FC-authorized recovery events. Saved ordinary deployment settings remain in the rocket and resume when FC simulation is disabled. There is no automatic fallback for an unassigned parachute.

The simulation's FC configuration maps shared logical channels 0--5 to recovery-device UUIDs and per-device delays. One channel can command several devices; several channels can provide redundant commands for one device. The bridge schedules only the first accepted request for each device, so repeated pulses cannot restart its delay.

\begin{equation}
 t_{\mathrm{deploy}}=t_{\mathrm{accepted\ command,boot}}-t_{\mathrm{warmup}}+\Delta t_{\mathrm{device}}.
\end{equation}
The command timestamp is the completion time of the relevant FC/Java action, retained at microsecond precision. A pyro pulse's duration is not a deployment delay. Warm-up commands before physical time zero are discarded physically, and commands after landing cannot deploy a device. Recovery timing does not wait for PWM or telemetry reception.

\section{Actual deployment and descent}
When handling a deployment event, the engine checks that the stage is active. It retains the existing under-thrust abort behavior, launch-rod warning, and high-speed-deployment warning. Listeners can veto events or deployment. A successful deployment adds the device to the status's deployed set, records the event, acknowledges FC provenance when applicable, and switches to the landing stepper if still airborne.

The landing stepper sums effective drag areas for deployed instances:
\begin{equation}
 C_{D,\mathrm{landing}}=\sum_{j\in\mathrm{deployed}}
 n_j C_{D,j}\frac{A_j}{S_{\mathrm{ref}}}.
\end{equation}
It then uses the descent integration path with atmospheric conditions, mass, gravity, and wind. The sum is the recovery-stepper drag model; it should not be assumed to be a complete addition of every ascent-airframe aerodynamic term. Ensemble drag scaling is explicitly applied to this descent coefficient too.

Drogue deployment does not end the FC. The same per-flight runtime continues while the landing stepper advances time, allowing a later main event. The bridge records actual device name/ID, origin channel, due time, deployment time, and lateness. Those events survive ORK storage and supported log import and drive the 3D canopy display.

\section{Stage separation and branch identity}
An ordinary staged simulation can split into multiple branches. At separation, the engine copies status, creates a booster result branch, clears the appropriate active-stage subsets in booster and continuing configurations, removes unattached pending events, and queues the detached branch for simulation. Each branch uses its remaining hardware and motor/recovery state.

The FC path currently requires exactly one active stage at startup and does not support arbitrary multistage FC ownership or mid-flight checkpoint boot. The general engine's ability to branch should not be mistaken for support for an independent FC on every resulting stage.

\section{Model boundaries and debugging order}
Opening inflation, reefing/cutaway, tether tension, separation impulses, electrical firing energy, and detailed canopy attitude dynamics are outside the current recovery model. Continuity/status bits are logical firmware observations, not evidence of a successful physical canopy deployment.

To investigate a missed deployment, follow the chain in order: logical output present; arming/accepted command; channel mapping; device active; warm-up/landing eligibility; scheduled event; listener acceptance; engine abort conditions; actual deployment record; descent drag. This identifies whether the problem is firmware decision logic, mapping, event authority, or physical integration without conflating those layers.

\sourceline{\source{core/src/main/java/info/openrocket/core/rocketcomponent/RecoveryDevice.java}{RecoveryDevice}; \source{core/src/main/java/info/openrocket/core/rocketcomponent/Parachute.java}{Parachute}; \source{core/src/main/java/info/openrocket/core/rocketcomponent/DeploymentConfiguration.java}{DeploymentConfiguration}; \source{core/src/main/java/info/openrocket/core/simulation/BasicLandingStepper.java}{BasicLandingStepper}; \source{core/src/main/java/info/openrocket/core/simulation/flightcomputer/FlightComputerRecoveryBridge.java}{FlightComputerRecoveryBridge}; \source{core/src/main/java/info/openrocket/core/simulation/ModifiedEventSimulationEngine.java}{ModifiedEventSimulationEngine}.}
